\SourceLeaf{192}
possunt, designante \(D\) determinantem formarum \(\Phi,\varphi\); \(m\) divisorem communem maximum numerorum \(a,2b,c\) (art.~162). Hoc igitur problema, quod pro valore negativo ipsius \(D\) iam supra solvimus, nunc pro positivo aggrediemur. Quia vero manifesto quivis valor ipsius \(t\) aequationi satisfaciens etiamnum mutato signo satisfacit, similiterque quivis valor ipsius \(u\): sufficiet si omnes valores \textit{positivos} ipsorum \(t,u\) assignare possimus, fungeturque quaelibet solutio per valores positivos, quatuor solutionum vice. Hoc negotium ita absolvemus, ut \textit{primo} valores \textit{minimos} ipsorum \(t,u\) (praeter hos per se obvios \(t=m,\ u=0\)) invenire, \textit{tum} ex his omnes reliquos derivare doceamus.

\Needspace{8\baselineskip}
\Article{198}
\textsc{Problema.} \textit{Invenire numeros minimos \(t,u\) aequationi indeterminatae \(tt-Duu=mm\) satisfacientes, siquidem forma aliqua \((M,N,P)\) datur, cuius determinans est \(D\), numerorumque \(M,2N,P\) divisor communis maximus \(n\).}

\textit{Sol.} Accipiatur ad lubitum forma reducta \((a,b,-a')\ldots f\), determinantis \(D\), ubi divisor communis maximus numerorum \(a,2b,a'\) sit \(m\), qualem dari vel inde manifestum est, quod forma reducta formae \((M,N,P)\) aequivalens inveniri potest, quae per art.~161 hac proprietate erit praedita: sed ad propositum praesens quaevis forma reducta in qua conditio haec locum habet poterit adhiberi. Evolvatur periodus formae \(f\), quam ex \(n\) formis constare supponemus. Retentis omnibus signis quibus in art.~188 usi sumus, \(f^n\) erit \((+a^n,b^n,-a^{n+1})\), quia \(n\) par, et in hanc formam transit \(f\) per substitutionem propriam \(\alpha^n,\beta^n,\gamma^n,\delta^n\). Quia vero \(f\) et \(f^n\) sunt identicae: \(f\) transit in \(f^n\) etiam per substitutionem propriam \(1,0,0,1\). Ex his duabus transformationibus similibus formae \(f\) in \(f^n\) per art.~162 deduci poterit solutio aequationis \(tt-Duu=mm\) \textit{in integris}, scilicet \(t=\frac12(\alpha^n+\delta^n)m\) (aequ.~18 art.~162), \(u=\frac{\gamma^n m}{a}\) (aequ.~19)\HistoricalNote{*)}{\begin{tabular}[t]{@{}l@{}}\(\begin{array}{@{}rll@{}}
\text{Quae in art. 162 erant}&\alpha,\beta,\gamma,\delta;&\alpha',\beta',\gamma',\delta';\\
\text{hic sunt}&1,0,0,1;&\alpha^n,\beta^n,\gamma^n,\delta^n;\end{array}\)
\\[0.6ex]\(\begin{array}{@{}lll@{}}A,B,C;&a,b,c;&e,\\a,b,-a';&a,b,-a';&1.\end{array}\)\end{tabular}}. Designentur hi valores positive accepti si forte nondum sunt per \(T,U\), eruntque hi \(T,U\) valores minimi ipsorum \(t,u\), praeter hos \(t=m,\ u=0\) (a quibus necessario erunt diversi. quia manifesto \(\gamma^n\) non poterit esse \(=0\)).

Supponamus enim dari adhuc minores valores ipsorum \(t,u\) puta \(\mathfrak{t},\mathfrak{u}\) qui sint positivi et \(\mathfrak{u}\) non \(=0\). Tum per art.~162 forma \(f\) per substitutionem propriam \(\frac1m(\mathfrak{t}-b\mathfrak{u}),\ \frac1m a'\mathfrak{u},\ \frac1m a\mathfrak{u},\ \frac1m(\mathfrak{t}+b\mathfrak{u})\) transformabitur in formam cum ipsa
\SourceLeaf{193}
identicam. Iam ex art.~193, II sequitur, \textit{aut} \(\frac1m(\mathfrak{t}-b\mathfrak{u})\) \textit{aut} \(-\frac1m(\mathfrak{t}-b\mathfrak{u})\) alicui numerorum \(\alpha'',\alpha''',\alpha''''\) etc. aequalem esse debere, puta \(=\alpha^\mu\) (quia enim \(\mathfrak{tt}=D\mathfrak{uu}+mm=bb\mathfrak{uu}+aa'\mathfrak{uu}+mm\), erit \(\mathfrak{tt}>bb\mathfrak{uu}\), adeoque \(\mathfrak{t}-b\mathfrak{u}\) positivus; hinc fractio \(\frac{\mathfrak{t}-b\mathfrak{u}}{a\mathfrak{u}}\), quae respondet fractioni \(\frac{\mathfrak{A}}{\mathfrak{C}}\) in art.~193, idem signum habebit ut \(a\) vel \(a'\)); atque in casu priori \(\frac1m a'\mathfrak{u},\ \frac1m a\mathfrak{u},\ \frac1m(\mathfrak{t}+b\mathfrak{u})\), in posteriori easdem quantitates mutatis signis, resp. \(=\beta^\mu,\gamma^\mu,\delta^\mu\). Sed quum sit \(\mathfrak{u}<U\) \textit{i.~e.} \(\mathfrak{u}<\frac{\gamma^n m}{a}\) et \(>0\): erit \(\gamma^\mu<\gamma^n\) et \(>0\); quocirca quum progressio \(\gamma,\gamma',\gamma''\) etc. continuo crescat, necessario \(\mu\) iacebit inter \(0\) et \(n\) excl. Forma vero respondens, \(f^\mu\), identica erit cum forma \(f\), \textit{Q.~E.~A.}, quum omnes formae \(f,f',f''\) etc. usque ad \(f^{n-1}\) diversae esse supponantur. Ex his colligitur, minimos valores ipsorum \(t,u\) (exceptis valoribus \(m,0\)) esse \(T,U\).

\textit{Ex.} Si \(D=79,\ m=1\): adhiberi poterit forma \((3,8,-5)\), pro qua \(n=6\), atque \(\alpha^n=-8,\ \gamma^n=-27,\ \delta^n=-152\) (art.~188). Hinc \(T=80,\ U=9\), qui sunt valores minimi numerorum \(t,u\), aequationi \(tt-79uu=1\) satisfacientes.

\Needspace{8\baselineskip}
\Article{199}
Ad praxin formulae adhuc commodiores erui possunt. Erit nimirum \(2b\gamma^n=-a(\alpha^n-\delta^n)\), quod facile ex art.~162 deducitur, multiplicando aequ.~\([19]\) per \(2b\), \([20]\) per \(a\) et mutando characteres illic adhibitos in praesentes. Hinc fit \(\alpha^n+\delta^n=2\delta^n-\frac{2b}{a}\gamma^n\), adeoque
\[
\pm T=m\left(\delta^n-\frac{b}{a}\gamma^n\right),\quad \pm U=\frac{\gamma^n m}{a}
\]
Per similem methodum hos valores obtinemus
\[
\pm T=m\left(\alpha^n+\frac{b}{a'}\beta^n\right),\quad \pm U=\frac{\beta^n m}{a'}
\]
Tum hae tum illae formulae perquam commodae evadunt, propter \(\gamma^n=\delta^{n-1}\), \(\alpha^n=\beta^{n-1}\), ita ut si his uteris, solam progressionem \(\beta',\beta'',\beta'''\ldots\beta^n\); si illis uti mavis, solam hanc \(\delta',\delta'',\delta'''\) etc. supputavisse sufficiat. Praeterea ex art.~189, 3 facile deducitur, quum \(n\) necessario sit par, \(\alpha^n\) et \(\frac{b}{a'}\beta^n\) eadem signa habere; neque minus \(\delta^n\) et \(\frac{b}{a}\gamma^n\), ita ut in formula priori pro \(T\) differentia absoluta, in posteriori summa absoluta accipi debeat, neque adeo ad signa respicere omnino opus sit. Receptis signis in art.~189, 4 adhibitis erit ex formula priori
\[
\begin{aligned}T&=m[k',k'',k'''\ldotp\,\ldotp\,\ldotp\,\ldotp\, k^n]-\frac{mb}{a}[k',k'',k'''\ldotp\,\ldotp\,\ldotp\,\ldotp\, k^{n-1}],\\U&=\frac{m}{a}[k',k'',k'''\ldotp\,\ldotp\,\ldotp\,\ldotp\, k^{n-1}]\end{aligned}
\]
\SourceLeaf{194}
ex posteriori
\[
\begin{aligned}T&=m[k'',k'''\ldotp\,\ldotp\,\ldotp\,\ldotp\, k^{n-1}]+\frac{mb}{a'}[k'',k'''\ldotp\,\ldotp\,\ldotp\,\ldotp\, k^n],\\U&=\frac{m}{a'}[k'',k'''\ldotp\,\ldotp\,\ldotp\,\ldotp\, k^n]\end{aligned}
\]
ubi pro valore ipsius \(T\) etiam \(m\left[k'',k'''\ldotp\,\ldotp\,\ldotp\,\ldotp\, k^n,\frac{b}{a'}\right]\) scribi poterit.

\textit{Ex.} Pro \(D=61,\ m=2\) adhiberi potest forma \((2,7,-6)\), pro qua eruitur \(n=6\); \(k',k'',k''',k'''',k''''',k''''''\) resp. \(=2,2,7,2,2,7\). Hinc fit
\[
\begin{aligned}T&=2[2,2,7,2,2,7]-7[2,2,7,2,2]\\&=2888-1365=1523\end{aligned}
\]
ex formula prima; idem provenit ex secunda
\[
T=2[2,7,2,2]+\tfrac73[2,7,2,2.7]
\]
\[
U\text{ vero fit }=[2,2,7,2,2]=\tfrac13[2,7,2,2,7]=195.
\]
Ceterum plura artificia adhuc dantur, per quae calculus contrahi potest, sed de his fusius hic loqui brevitas non permittit.

\Needspace{8\baselineskip}
\Article{200}
Ut ex valoribus minimis ipsorum \(t,u\) \textit{omnes} obtineamus, aequationem \(TT-DUU=mm\) ita exhibemus
\[
\left(\frac{T}{m}+\frac{U}{m}\surd D\right)\left(\frac{T}{m}-\frac{U}{m}\surd D\right)=1
\]
unde etiam erit
\[
\left(\frac{T}{m}+\frac{U}{m}\surd D\right)^e\left(\frac{T}{m}-\frac{U}{m}\surd D\right)^e=1\quad\mathinner{\ldotp\ldotp\ldotp\ldotp\ldotp\ldotp\ldotp}\tag*{[1]}
\]
denotante \(e\) numerum quemcunque. Iam designabimus brevitatis caussa valores quantitatum
\begingroup\postdisplaypenalty=10000
\[
\begin{gathered}\frac{m}{2}\left(\frac{T}{m}+\frac{U}{m}\surd D\right)^e+\frac{m}{2}\left(\frac{T}{m}-\frac{U}{m}\surd D\right)^e,\\
\frac{m}{2\surd D}\left(\frac{T}{m}+\frac{U}{m}\surd D\right)^e-\frac{m}{2\surd D}\left(\frac{T}{m}-\frac{U}{m}\surd D\right)^e\text{\textsuperscript{*)}}\end{gathered}
\]
\HistoricalDisplayNote{*)}{In his solis quatuor expressionibus et in aequ.~\([1]\) \(e\) denotat \textit{exponentem} potestatis: in reliquis literae apici adscriptae semper \textit{indicem} designant.}\endgroup
generaliter per \(t^e,u^e\) resp. \textit{i.~e.} illarum valores pro \(e=0\), per \(t^0,u^0\) (qui erunt \(m,0\)); pro \(e=1\) per \(t',u'\) (qui fiunt \(T,U\)); pro \(e=2\) per \(t'',u''\); pro \(t=3\) per \(t''',u'''\) etc. --- demonstrabimusque, si pro \(e\) accipiantur omnes numeri integri non negativi \textit{i.~e.} \(0\), omnesque positivi ab \(1\) usque ad \(\infty\), expressiones illas exhibere omnes valores positivos ipsorum \(t,u\): scilicet I) omnes valores illarum expressio%
\SourceLeaf{195}%
num esse revera valores ipsorum \(t,u\); II) omnes valores illos esse numeros integros; III) nullos valores positivos ipsorum \(t,u\) dari qui sub formulis illis non contineantur.

I. Substitutis pro \(t^e,u^e\) valoribus suis nullo negotio adiumento aequ.~\([1]\) confirmatur, esse
\[
(t^e+u^e\surd D)(t^e-u^e\surd D)=mm,\quad \textit{i.~e.}\quad t^e t^e-Du^e u^e=mm
\]

II. Eodem modo facile confirmatur, esse generaliter
\[
t^{e+1}+t^{e-1}=\frac{2T}{m}t^e,\quad u^{e+1}+u^{e-1}=\frac{2T}{m}u^e
\]
Hinc manifestum est, duas progressiones \(t^0,t',t'',t'''\) etc., \(u^0,u',u'',u'''\) etc. esse recurrentes, et utriusque scalam relationis \(\frac{2T}{m},-1\), scilicet
\[
t''=\frac{2T}{m}t'-t^0,\quad t'''=\frac{2T}{m}t''-t'\text{ etc.,}\quad u''=\frac{2T}{m}u'\text{ etc.}
\]
Iam quoniam per hyp. forma aliqua datur, \((M,N,P)\), determinantis \(D\), in qua \(M,2N,P\) per \(m\) sunt divisibiles: habebitur
\[
TT=(NN-MP)UU+mm
\]
eritque adeo manifesto \(4TT\) per \(mm\) divisibilis. Hinc \(\frac{2T}{m}\) erit numerus integer et quidem positivus. Quia vero \(t^0=m,\ t'=T,\ u^0=0,\ u'=U\), adeoque integri: omnes \(t'',t'''\) etc. \(u'',u'''\) etc. etiam integri erunt. Porro perspicuum est, quia \(TT>mm\), omnes \(t^0,t',t'',t'''\) etc. positivos et continuo in infinitum crescentes esse, nec non omnes \(u^0,u',u'',u'''\) etc.

III. Supponamus, dari adhuc alios valores positivos ipsorum \(t,u\) qui in progressione \(t^0,t',t''\) etc. \(u^0,u',u''\) etc. non contenti sint, puta \(\mathfrak{T},\mathfrak{U}\). Manifestum est. quum progressio \(u^0,u'\) etc. a \(0\) in infinitum crescat, \(\mathfrak{U}\) necessario inter duos terminos proximos, \(u^n\) et \(u^{n+1}\) situm fore, ita ut sit \(\mathfrak{U}>u^n\) et \(\mathfrak{U}<u^{n+1}\). Ut absurditatem, huius suppositionis demonstremus, observamus

\(1^\circ\) Aequationi \(tt-Duu=mm\) satisfactum iri etiam ponendo
\[
t=\frac1m(\mathfrak{T}t^n-D\mathfrak{U}u^n),\quad u=\frac1m(\mathfrak{U}t^n-\mathfrak{T}u^n).
\]
\SourceLeaf{196}
Hoc quidem nullo negotio per substitutionem confirmatur: quod vero hi valores quos ponemus brevitatis gratia \(=\tau,\upsilon\), semper sunt numeri \textit{integri}, ita ostendimus. Si \((M,N,P)\) est forma determinantis \(D\), atque \(m\) divisor communis numerorum \(M,2N,P\): erit tum \(\mathfrak{T}+N\mathfrak{U}\) tum \(t^n+Nu^n\) per \(m\) divisibilis adeoque etiam \(\mathfrak{U}(t^n+Nu^n)-u^n(\mathfrak{T}+N\mathfrak{U})\) sive \(\mathfrak{U}t^n-\mathfrak{T}u^n\). Quare \(\upsilon\) erit integer et proin etiam \(\tau\), quia \(\tau\tau=D\upsilon\upsilon+mm\).

\(2^\circ\) Patet \(\upsilon\) non posse esse \(=0\); hinc enim sequeretur
\[
\mathfrak{UU}t^n t^n=\mathfrak{TT}u^n u^n
\]
sive
\[
\mathfrak{UU}(Du^n u^n+mm)=u^n u^n(D\mathfrak{UU}+mm)
\]
sive \(\mathfrak{UU}=u^n u^n\), contra hyp. ex qua \(\mathfrak{U}>u^n\). Quum igitur praeter valorem \(0\), minimus valor ipsius \(u\) sit \(U\), erit \(\upsilon\) certe non minor quam \(U\).

\(3^\circ\) Facile ex valoribus ipsorum \(t^n,t^{n+1},u^n,u^{n+1}\) confirmari potest, esse
\[
mU=u^{n+1}t^n-t^{n+1}u^n.
\]
Quare \(\mathfrak{U}t^n-\mathfrak{T}u^n\) certe non erit minor quam \(u^{n+1}t^n-t^{n+1}u^n\).

\(4^\circ\) Iam ex aequatione \(\mathfrak{TT}-D\mathfrak{UU}=mm\) habetur
\[
\frac{\mathfrak{T}}{\mathfrak{U}}=\surd\left(D+\frac{mm}{\mathfrak{UU}}\right)
\]
et similiter
\[
\frac{t^{n+1}}{u^{n+1}}=\surd\left(D+\frac{mm}{u^{n+1}u^{n+1}}\right)
\]
unde facile deducitur esse \(\frac{\mathfrak{T}}{\mathfrak{U}}>\frac{t^{n+1}}{u^{n+1}}\). Hinc vero et ex conclusione in \(3^\circ\) sequitur
\[
\begin{aligned}&(\mathfrak{U}t^n-\mathfrak{T}u^n)\left(t^n+u^n\frac{\mathfrak{T}}{\mathfrak{U}}\right)\\&\qquad>(u^{n+1}t^n-t^{n+1}u^n)\left(t^n+u^n\frac{t^{n+1}}{u^{n+1}}\right)\end{aligned}
\]
sive, evolutione facta, et loco ipsorum \(\mathfrak{TT},\ t^n t^n,\ t^{n+1}t^{n+1}\) substitutis valoribus suis \(D\mathfrak{UU}+mm,\ Du^n u^n+mm,\ Du^{n+1}u^{n+1}+mm\),
\[
\frac1{\mathfrak{U}}(\mathfrak{UU}-u^n u^n)>\frac1{u^{n+1}}(u^{n+1}u^{n+1}-u^n u^n)
\]
unde, quoniam utraque quantitas manifesto positiva, fit transponendo \(\mathfrak{U}+\frac{u^n u^n}{u^{n+1}}>u^{n+1}+\frac{u^n u^n}{\mathfrak{U}}\), \textit{Q.~E.~A.}, quia quantitatis prioris pars prima \textit{minor} est
\SourceLeaf{197}
quam pars prima quantitatis secundae, nec non illius secunda minor quam secunda huius. Quamobrem suppositio consistere nequit et progressiones \(t^0,t',t''\) etc. \(u^0,u',u''\) etc. omnes valores positivos ipsorum \(t,u\) exhibebunt.

\textit{Ex.} Pro \(D=61,\ m=2\) valores minimos positivos ipsorum \(t,u\) invenimus \(1523,195\): quare \textit{omnes} valores positivi exhibebuntur per has formulas
\[
\begin{aligned}
t&=\left(\frac{1523}{2}+\frac{195}{2}\surd61\right)^e+\left(\frac{1523}{2}-\frac{195}{2}\surd61\right)^e\\
u&=\frac1{\surd61}\left(\left(\frac{1523}{2}+\frac{195}{2}\surd61\right)^e-\left(\frac{1523}{2}-\frac{195}{2}\surd61\right)^e\right)
\end{aligned}
\]
Invenitur autem
\[
\begin{aligned}
t^0&=2,\quad t'=1523,\\
t''&=1523t'-t^0=2319527,\\
t'''&=1523t''-t'=3532618098\text{ etc.}\\
u^0&=0,\quad u'=195,\\
u''&=1523u'-u^0=296985,\\
u'''&=1523u''-u'=452307960\text{ etc.}
\end{aligned}
\]

\Needspace{8\baselineskip}
\Article{201}
Circa problema in artt. praecc. tractatum sequentes observationes adhuc adiicimus.

1) Quum aequationem \(tt-Duu=mm\) pro omnibus casibus solvere docuerimus, ubi \(m\) est divisor communis maximus trium numerorum \(M,2N,P\), talium ut \(NN-MP=D\): operae pretium est omnes numeros qui tales divisores esse possunt sive omnes valores ipsius \(m\) pro valore dato ipsius \(D\) assignare. Ponatur \(D=nnD'\), ita ut \(D'\) a factoribus quadraticis omnino sit liber, quod obtinetur si pro \(nn\) assumitur maximum quadratum ipsum \(D\) metiens: sin vero \(D\) iam per se nullum factorem quadraticum implicaret, fieri deberet \(n=1\). Tum dico

\textit{Primo}, si \(D'\) fuerit formae \(4k+1\), quemvis divisorem ipsius \(2n\) fore valorem ipsius \(m\), et vice versa. Si enim \(g\) est divisor ipsius \(2n\), habebitur forma \(\left(g,n,\frac{nn(1-D')}{g}\right)\), cuius determinans est \(D\), et in qua manifesto divisor communis maximus numerorum \(g,2n,\frac{nn(D'-1)}{g}\) erit \(g\) (patet enim \(\frac{nn(D'-1)}{gg}=\frac{4nn}{gg}\cdot\frac{D'-1}{4}\) esse numerum integrum). Si vero, vice versa, \(g\) supponitur esse valor ipsius \(m\), scilicet divisor communis maximus numerorum \(M,2N,P\), atque \(NN-MP=D\): manifesto \(4D\) sive \(4nnD'\) divisibilis erit per \(gg\). Hinc vero sequitur, \(2n\) necessario per \(g\) divisibilem esse. Si enim \(g\) ipsum \(2n\) non metiretur, \(g\) et \(2n\) haberent divisorem communem maximum minorem quam \(g\), quo posito \(=\delta\), at%
\SourceLeaf{198}%
que \(2n=\delta n',\ g=\delta g'\), foret \(n'n'D'\) per \(g'g'\) divisibilis, \(n'\) ad \(g'\) adeoque etiam \(n'n'\) ad \(g'g'\) primus et proin etiam \(D'\) per \(g'g'\) divisibilis, contra hyp. secundum quam \(D'\) ab omni factore quadratico est liberatus.

\textit{Secundo}, si \(D'\) fuerit formae \(4k+2\) vel \(4k+3\). quemvis divisorem ipsius \(n\) fore valorem ipsius \(m\), et vice versa quemvis valorem ipsius \(m\) metiri ipsum \(n\). Si enim \(g\) est divisor ipsius \(n\), habebitur forma \(\left(g,0,-\frac{nnD'}{g}\right)\), cuius determinans \(=D\), et ubi manifesto numerorum \(g,0,\frac{nnD'}{g}\) divisor communis maximus erit \(g\). --- Si vero \(g\) supponitur esse valor ipsius \(m\), puta divisor communis maximus numerorum \(M,2N,P\), atque \(NN-MP=D\): eodem modo ut supra \(g\) metietur ipsum \(2n\), sive \(\frac{2n}{g}\) erit integer. Si quotiens hic esset impar: quadratum \(\frac{4nn}{gg}\) foret \(\equiv1\ (\mathrm{mod.}\ 4)\), adeoque \(\frac{4nnD'}{gg}\) aut \(\equiv2\) aut \(\equiv3\ (\mathrm{mod.}\ 4)\). At \(\frac{4nnD'}{gg}=\frac{4D}{gg}=\frac{4NN}{gg}-\frac{4MP}{gg}\equiv\frac{4NN}{gg}(\mathrm{mod.}\ 4)\), et proin \(\frac{4NN}{gg}\) aut \(\equiv2\) aut \(\equiv3\ (\mathrm{mod.}\ 4)\), \textit{Q.~E.~A.}, quia omne quadratum aut cifrae aut unitati secundum modulum \(4\) congruum esse debet. Quare quotiens \(\frac{2n}{g}\) necessario erit par, adeoque \(\frac{n}{g}\) integer, sive \(g\) divisor ipsius \(n\).

Patet itaque, \(1\) semper esse valorem ipsius \(m\), sive aequationem \(tt-Duu=1\) pro quovis valore positivo non quadrato ipsius \(D\) per praecedentia resolubilem esse; \(2\) tunc tantummodo esse valorem ipsius \(m\), si \(D\) fuerit aut formae \(4k\), aut formae \(4k+1\).

2) Si \(m\) est maior quam \(2\), attamen numerus idoneus, solutio aequationis \(tt-Duu=mm\) reduci potest ad solutionem similis aequationis, ubi \(m\) est aut \(1\) aut \(2\). Scilicet posito ut ante \(D=nnD'\), si \(m\) ipsum \(n\) metitur, metietur \(mm\) ipsum \(D\). Tum si valores minimi ipsorum \(p,q\) in aequatione \(pp-\frac{D}{mm}qq=1\) supponuntur esse \(p=P,\ q=Q\), valores minimi ipsorum \(t,u\) in aequatione \(tt-Duu=mm\) erunt \(t=mP,\ u=Q\). --- Si vero \(m\) ipsum \(n\) non metitur, metietur saltem ipsum \(2n\) eritque certo par; \(\frac{4D}{mm}\) autem integer. Et si tunc valores minimi ipsorum \(p,q\) in aequatione \(pp-\frac{4D}{mm}qq=4\) inventi sunt \(p=P,\ q=Q\): valores minimi ipsorum \(t,u\) in aequatione \(tt-Duu=mm\) erunt \(t=\frac{m}{2}P,\ u=Q\). --- In utroque autem casu non solum ex valoribus minimis ipsorum \(p,q\) valores minimi ipsorum \(t,u\), sed ex \textit{omnibus} valoribus illorum \textit{omnes} valores horum per hanc methodum manifesto deduci poterunt.

3) Designantibus \(t^0,u^0;\ t',u';\ t'',u''\) etc. omnes valores positivos ipsorum
\SourceLeaf{199}
\(t,u\) in aequatione \(tt-Duu=mm\) (ut in art. praec.), si contingit ut valores quidam ex serie illa, valoribus primis in eadem secundum modulum quemcunque datum \(r\), congrui sint, puta \(t^\rho\equiv t\) (sive \(\equiv m\)), \(u^\rho\equiv u^0\) sive \(\equiv0\ (\mathrm{mod.}\ r)\); simulque valores proxime sequentes valoribus secundis, puta
\[
t^{\rho+1}\equiv t',\quad u^{\rho+1}\equiv u'\quad(\mathrm{mod.}\ r)
\]
erit etiam
\[
t^{\rho+2}\equiv t'',\quad u^{\rho+2}\equiv u'';\quad t^{\rho+3}\equiv t''',\quad u^{\rho+3}\equiv u'''\text{ etc.}
\]
Hoc facile inde deducitur, quod utraque series \(t^0,t',t''\) etc., \(u^0,u',u''\) etc. est ex recurrentium genere, scilicet quoniam
\[
t''=\frac{2T}{m}t'-t^0,\quad t^{\rho+2}=\frac{2T}{m}t^{\rho+1}-t^\rho
\]
erit
\[
t''\equiv t^{\rho+2}
\]
similiterque de reliquis. --- Hinc autem sequitur, fore generaliter
\[
t^{h+\rho}\equiv t^h,\quad u^{h+\rho}\equiv u^h(\mathrm{mod.}\ r)
\]
denotante \(h\) numerum quemcunque, nec non adhuc generalius, si fuerit
\[
\mu\equiv\nu\ (\mathrm{mod.}\ \rho),\quad\text{fore }t^\mu\equiv t^\nu,\quad u^\mu\equiv u^\nu(\mathrm{mod.}\ r).
\]

4) Conditionibus autem in observ. praec. requisitis semper satisfieri potest, scilicet semper inveniri potest index \(\rho\) (pro modulo quocunque dato \(r\)), pro quo sit
\[
t^\rho\equiv t^0,\quad t^{\rho+1}\equiv t',\quad u^\rho\equiv u^0,\quad u^{\rho+1}\equiv u'
\]
Ad quod demonstrandum observamus

\textit{primo}, conditioni tertiae semper satisfieri posse. Nullo enim negotio per criteria in (1) tradita perspicietur, etiam aequationem \(pp-rrDqq=mm\) solubilem fore; et si valores minimi positivi ipsorum \(p,q\) (praeter hos \(m,0\)) supponuntur esse \(P,Q\): inter valores ipsorum \(t,u\) manifesto erunt etiam \(t=P,\ u=rQ\). Quare \(P,rQ\) in progressionibus \(t^0,t'\) etc., \(u^0,u'\) etc. contenti erunt, et si \(P=t^\lambda,\ rQ=u^\lambda\), erit \(u^\lambda\equiv0\equiv u^0(\mathrm{mod.}\ r)\). Praeterea facile perspicietur, inter \(u^0\) et \(u^\lambda\) nullum terminum fore ipsi \(u^0\) secundum modulum \(r\) congruum.

\textit{Secundo} patet, si hic insuper tres reliquae conditiones adimpletae sint, puta si etiam \(u^{\lambda+1}\equiv u',\ t^\lambda\equiv t^0,\ t^{\lambda+1}\equiv t'\), poni tantummodo debere \(\rho=\lambda\). Si
\SourceLeaf{200}
vero una aut altera illarum conditionum locum non habet, dico certe statui posse \(\rho=2\lambda\). Nam ex aequat.~\([1]\) formulisque generalibus pro \(t^e,u^e\) in art.~praec. deducitur
\[
t^{2\lambda}=\frac1m(t^\lambda t^\lambda+Du^\lambda u^\lambda)=\frac1m(mm+2Du^\lambda u^\lambda)
\]
adeoque
\[
\frac{t^{2\lambda}-t^0}{r}=\frac{2Du^\lambda u^\lambda}{mr}
\]
quae quantitas erit numerus integer, quia per hyp. \(r\) ipsum \(u^\lambda\) metitur, nec non \(mm\) ipsum \(4D\), adeoque a potiori \(m\) ipsum \(2D\). --- Porro erit \(u^{2\lambda}=\frac2m t^\lambda u^\lambda\), et quoniam
\[
4t^\lambda t^\lambda=4Du^\lambda u^\lambda+4mm
\]
adeoque per \(mm\) divisibilis, \(2t^\lambda\) erit divisibilis per \(m\), et proin \(u^{2\lambda}\) per \(r\), sive
\[
u^{2\lambda}\equiv u^0\ (\mathrm{mod.}\ r)
\]
Tertio invenitur
\[
t^{2\lambda+1}=t'+\frac{2Du^\lambda u^{\lambda+1}}{m}
\]
et quoniam ex simili ratione \(\frac{2Du^\lambda}{mr}\) est integer, erit
\[
t^{2\lambda+1}\equiv t'\ (\mathrm{mod.}\ r)
\]
Tandem reperitur
\[
u^{2\lambda+1}=u'+\frac{2t^{\lambda+1}u^\lambda}{m}
\]
et quoniam \(2t^{\lambda+1}\) per \(m\) divisibilis est, \(u^\lambda\) per \(r\): erit
\[
u^{2\lambda+1}\equiv u'\ (\mathrm{mod.}\ r).\quad \textit{Q.~E.~D.}
\]

Ceterum usus posteriorum duarum observationum in sequentibus apparebit.

\Needspace{8\baselineskip}
\Article{202}
Casus particularis problematis, nempe solvere aequationem \(tt-Duu=1\), iam a geometris seculi praecedentis fuit agitatus. Sagacissimus Fermatius problema hoc analystis Anglis proposuit, Wallisiusque Brounkerum tamquam inventorem solutionis, quam in \textit{Alg. Cap.}~98, \textit{Opp. T.}~II \textit{p.}~418 sqq. tradit, nominat; Ozanam Fermatium; denique ill. Euler, qui de illo egit in \textit{Comm. Petr.}~VI \textit{p.}~175, \textit{Comm. nov.}~XI \textit{p.}~28\HistoricalNote{*)}{In hac comm. algorithmus quem art.~27 exposuimus, per similia signa exhibetur, quod nos illic annotare negleximus.}, \textit{Algebra P.}~II \textit{p.}~226, \textit{Opusc. An.}~I \textit{p.}~310, Pellium. unde
\SourceLeaf{201}
problema illud a quibusdam auctoribus \textit{Pellianum} vocatum est. Omnes hae solutiones, si essentiam spectas, conveniunt cum ea quam obtinemus, si in art.~198 formam reductam eam adoptamus in qua \(a=1\); attamen operationem quam praescribunt tandem necessario \textit{finiri}, sive problema semper \textit{revera solubile} esse, nemo ante ill. La Grange rigorose\HistoricalNote{*)}{Quae Wallisius ad hunc finem protulit l.~c.~p.~427, 428 nihil ponderis habent. Paralogismus in eo consistit, quod p.~428 l.~4, supponit, proposita quantitate \(p\) inveniri posse numeros integros \(a,z\) tales ut \(\frac{z}{a}\) minor sit quam \(p\), defectus vero \textit{assignato} minor. Hoc utique verum est, quando defectus \textit{assignatus} est \textit{quantitas data}, neque vero, quando ab \(a\) et \(z\) pendet adeoque variabilis est, uti in casu praesenti evenit.} demonstravit, \textit{Mélanges de la Soc. de Turin T.}~IV \textit{p.}~19, et concinnius \textit{Hist. de l'Ac. de Berlin}, 1767, \textit{p.}~237. Exstat haec disquisitio etiam in \textit{supplementis ad Euleri Algebram} iam saepius laudatis. Ceterum methodus nostra (ex principiis omnino diversis petita, neque ad casum \(m=1\) restricta) plerumque plures vias ad solutionem perveniendi suppeditat, quoniam in art.~198 a quavis alia forma reducta \((a,b,-a')\) proficisci possumus.

\Needspace{8\baselineskip}
\Article{203}
\textsc{Problema.} \textit{Si formae \(\Phi,\varphi\) sunt aequivalentes, omnes transformationes alterius in alteram exhibere.}

\textit{Sol.} Quando formae hae unico tantum modo aequivalentes sunt (\textit{i.~e.} aut proprie tantum aut improprie tantum) quaeratur per art.~196 transformatio una formae \(\varphi\) in \(\Phi\), quae sit \(\alpha,\beta,\gamma,\delta\), patetque alias quam quae huic sint similes, dari non posse. Quando vero \(\varphi,\Phi\) tum proprie tum improprie aequivalent, quaerantur duae transformationes dissimiles. \textit{i.~e.} altera propria altera impropria, puta \(\alpha,\beta,\gamma,\delta\) et \(\alpha',\beta',\gamma',\delta'\), eritque quaevis alia transformatio aut huic aut illi similis. Si itaque forma \(\varphi\) est \((a,b,c)\), ipsius determinans \(=D\), divisor communis maximus numerorum \(a,2b,c\) (uti semper in praec.) \(m\), atque \(t,u\) indefinite omnes numeri aequationi \(tt-Duu=mm\) satisfacientes: in casu priori omnes transformationes formae \(\varphi\) in \(\Phi\) contentae erunt sub prima formularum sequentium I, in posteriori vel sub prima I vel sub secunda II.
\[
\begin{array}{@{}lll@{}}
\mathrm{I}\;\ldotp\,\ldotp\,\ldotp\,\ldotp\,\ldotp\,\ldotp\,\ldotp&\frac1m(\alpha t-(\alpha b+\gamma c)u),&\frac1m(\beta t-(\beta b+\delta c)u)\\[1ex]
&\frac1m(\gamma t+(\alpha a+\gamma b)u),&\frac1m(\delta t+(\beta a+\delta b)u)\\[1.5ex]
\mathrm{II}\;\ldotp\,\ldotp\,\ldotp\,\ldotp\,\ldotp\,\ldotp\,\ldotp&\frac1m(\alpha't-(\alpha'b+\gamma'c)u),&\frac1m(\beta't-(\beta'b+\delta'c)u)\\[1ex]
&\frac1m(\gamma't+(\alpha'a+\gamma'b)u),&\frac1m(\delta't+(\beta'a+\delta'b)u)
\end{array}
\]
